% ==============================================================
\section{はじめに}\label{sec:intro}
% ==============================================================

頻出パターンマイニング~\cite{agrawal1994apriori,han2000mining}は
トランザクションデータ中に繰り返し出現するアイテムセットを発見する基盤技術である。
近年、パターンの出現頻度（サポート）が\emph{なぜ}変動するかを
外部イベントに帰属させるフレームワークが提案されている~\cite{dsaa2025}。
このイベント帰属パイプラインは、サポート時系列上の変化点を検出し、
変化点とイベントの時間的近接度・変化量に基づくスコアリング、
円形シフト置換検定、BH補正、Union-Find重複排除の5ステップで構成される。

しかし、既存のイベント帰属パイプラインは\textbf{時間軸のみ}を扱う。
現実のデータは空間的位置（店舗、地域、国）などの補助次元を持つことが多く、
イベントの効果は空間的に局在する。例えば：

\begin{itemize}
\item 小売チェーンの夏季キャンペーンは沿岸部店舗でのみ売上を押し上げる。
\item 新規制の導入は特定の州・県でのみ処方パターンを変化させる。
\item オンラインプロモーションは特定の国・地域からのアクセスパターンに影響する。
\end{itemize}

1次元パイプラインでこれらを扱うと、2つの問題が生じる。
第一に、全空間にわたってサポートを集約するため、
局所的なイベント効果が\textbf{希釈}されて検出力が低下する。
第二に、イベントの空間的影響範囲が不明なまま帰属が行われるため、
「どの地域で効果があったか」という問いに答えられない。

これらの問題を解決するには、
サポートの時間的変動だけでなく空間的分布を同時にモデル化し、
イベントの時空間的スコープとの関連を検定する必要がある。
しかし、多次元化は以下の技術的課題を伴う：

\begin{enumerate}
\item \textbf{空間自己相関}：近接する空間位置のサポート時系列は正に相関する。
1次元パイプラインの円形シフト置換検定は空間構造を無視するため、
空間的にクラスター化した変化点に対して偽陽性が増加する。

\item \textbf{仮説空間の増大}：パターン$\times$イベントに空間次元が加わり、
多重検定の補正がより厳格になる必要がある。

\item \textbf{信号の空間的局在}：イベント効果が一部の空間領域に限られる場合、
空間全体の集約サポートでは信号対雑音比が低下する。
\end{enumerate}

\subsection{本研究の貢献}

関連研究の分析（\ref{sec:related}節）が示すように、
アイテムセットパターンのサポート変動を空間次元付きでイベントに帰属させる
フレームワークは存在しない（表~\ref{tab:related}）。
本稿ではこのギャップを埋める
\textbf{時空間イベント帰属パイプライン}を提案する。
貢献は以下の3点である。

\begin{enumerate}
\item \textbf{時空間帰属スコア}（\ref{sec:problem}節）：
サポート曲面上の各空間位置で変化点を検出し、
時間的近接度 $\mathrm{prox}_t$・空間的近接度 $\mathrm{prox}_s$・変化量 $\mathrm{mag}$ の
3要素から帰属スコアを定義する。
空間スキャン統計~\cite{kulldorff1997spatial}が空間クラスターの検出に留まるのに対し、
本スコアは検出された変動を既知イベントに帰属させる。

\item \textbf{時空間置換検定}（\ref{sec:method}節）：
サロゲートデータ法~\cite{lancaster2018surrogate}の時間的円形シフトを
空間構造保存に拡張した置換方式を設計する。
これにより空間自己相関構造を保ったまま帰無分布を生成し、
第一種過誤率を正しく制御する。

\item \textbf{包括的な実験評価}（\ref{sec:experiments}節）：
1次元パイプライン~\cite{dsaa2025}との比較を通じて、
(i) 空間的に局在するイベント効果の検出力向上、
(ii) イベント非対応パターンの偽帰属率低減を実証する。
\end{enumerate}
